The $3.65 T residual baseline is not a theoretical construct invented inside the model.  

It is the direct consequence of real-world money and claim issuance that has already occurred and continues to occur:

The actual U.S. federal deficit projected by the CBO for FY 2026 (≈ $1.90 T)
Documented global structural deficits, stock-flow adjustments, and off-balance-sheet claim expansions (≈ $1.75 T)

These are real fiscal and monetary events.  

Because true repayment rates (\(\rho_t\)) have historically been low, the large majority of these issued claims remain unextinguished. The resulting residual stock is therefore a factual consequence of real-life money and credit creation, not an abstract modeling choice.

The subsequent timeline (2055 median, 2059 full saturation) and the liquidation magnitudes follow as the mechanical result of that real residual once the locked dynamics are applied.
